\begin{abstract}
    \textit{Headline Incongruity Detection} (HID) adalah tugas mengidentifikasi ketidaksesuaian antara judul berita dan isi teksnya. Tugas ini menjadi kritis karena pembaca daring cenderung hanya memindai judul sehingga informasi yang keliru dapat menyebar luas, terutama pada isu sensitif. Tantangan utama pada dataset berita Indonesia adalah distribusi kelas yang sangat tidak seimbang (rasio 9:1), kemiripan leksikal tinggi di kedua kelas (rerata \textit{overlap} 0,796 versus 0,749), serta larangan penggunaan model \textit{pretrained}. Penelitian ini mengusulkan HIDRA (\textit{Headline Incongruity Detection with Rival-paired Monotonic Alignment}), suatu \textit{pipeline} yang menggabungkan rekayasa fitur multiblok dengan pemodelan paralel. Rekayasa fitur mencakup penyelarasan monoton berbasis \textit{longest increasing subsequence}, pemeringkatan rival berbasis indeks TF-IDF, serta perbandingan berpasangan terhadap tetangga terdekat. Pemodelan paralel terdiri dari \textit{ensemble} pohon bergradien beranggotakan enam model untuk menangkap interaksi struktural, serta \textit{cross encoder} BiGRU dengan mekanisme atensi yang dilatih pada data sintetis tujuh skema korupsi untuk menangkap interaksi kontekstual. Seluruh komponen dilatih dari awal tanpa bobot \textit{pretrained} dan divalidasi menggunakan \textit{StratifiedGroupKFold} lima lipatan. Evaluasi terhadap data validasi menghasilkan \textit{Macro F1} sebesar 0,8728, dengan skor kelas minoritas (\texttt{Tidak Sesuai}) mencapai 0,7690. Kontribusi utama meliputi fitur penyelarasan monoton dan perbandingan rival yang mengubah pertanyaan absolut menjadi komparatif, serta generator korupsi sintetis yang mengarahkan \textit{cross encoder} pada tipologi kesalahan nyata. Pengembangan lebih lanjut meliputi eksplorasi \textit{cross-lingual transfer} dan optimasi inferensi untuk penerapan waktu nyata.
\end{abstract}

\noindent\textbf{Kata Kunci:} \textit{Headline Incongruity Detection}, penyelarasan monoton, perbandingan rival, \textit{gradient boosted decision tree}, \textit{cross encoder}, deteksi ketidaksesuaian berita
